\chapter{Metodologia}
%TODO: PADRONIZAR NLP OU NLP ENTRE OUTRAS SIGLAS

\section{Processamento de linguagem natural e a evolução da tradução automática}

Neste capítulo, detalharemos a metodologia utilizada para desenvolver e avaliar um sistema de tradução automática (MT) voltado para a tradução jurídica, considerando as nuances culturais e legais. Utilizando Processamento de Linguagem Natural (NLP) e técnicas de aprendizado profundo, buscamos abordar as dificuldades específicas da tradução jurídica com o apoio de avaliações realizadas por humanos para refinar o sistema.

A metodologia adotada visa minimizar as limitações comuns nas traduções jurídicas automatizadas, como a interpretação incorreta de termos ambíguos e a falta de compreensão do contexto cultural e legal. Nosso objetivo é criar um modelo de NLP que não apenas traduza palavras, mas que também compreenda o significado subjacente dos termos e suas implicações jurídicas.

\section{Problemáticas clássicas de tradução jurídica}

Para enfrentar os desafios específicos da tradução jurídica, estrutura-se a problemática em dois tipos principais de dificuldades de tradução:

\subsection{Termos com significados diferentes em jurisdições distintas}

Esse problema ocorre quando dois sistemas legais utilizam palavras com conotações semânticas semelhantes, mas que entre duas línguas diferentes, possuem estrutura sintática completamente diferente, ainda mais fora de contexto. 
Por exemplo, no Brasil, o termo \textit{calção} refere-se a um valor monetário que assegura o cumprimento de um contrato, geralmente em contexto imobiliário. Já nos Estados Unidos, a palavra \textit{consideration} cumpre uma função similar, mas possui implicações próprias do direito norte-americano, como a garantia de um contrato por meio de um penhor de valor. Este estudo busca treinar o modelo para identificar tais termos, com auxilio de uma etiquetação ontológica do meio em que se identifica essas palavras, para refletir adequadamente as nuances de cada contexto legal.

%Evidenciar isso utilizando uma imagem _> Calção ->Calças/penhor
%Consideration -> consideração, penhor

\subsection{Diferenças conceituais para termos iguais}

Neste caso, a problemática se refere a termos idênticos que possuem interpretações e limitações legais variadas entre jurisdições. Um exemplo claro é o conceito de \textit{liberdade de expressão}, que na Constituição Americana é considerado um direito absoluto garantido pela Primeira Emenda. Em contraste, no Brasil, esse direito é defendido até o ponto em que não infrinja a liberdade de terceiros, o que implica uma limitação na interpretação desse termo. Considerando tais diferenças, a metodologia proposta inclui um mapeamento dessas limitações legais para que o sistema possa reconhecer estes termos, e \textbf{adicionar} contexto para refletir as nuanças do assunto.

\section{Tecnologia e ferramentas utilizadas}

Para desenvolver o sistema de tradução automática, utilizamos a linguagem de programação Python em conjunto com bibliotecas de aprendizado profundo, como \textit{TensorFlow} e \textit{PyTorch}, que facilitam o treinamento de modelos complexos de NLP. Essas bibliotecas permitem a manipulação de grandes quantidades de dados textuais e o desenvolvimento de redes neurais capazes de identificar padrões de linguagem e contexto.

Para a criação e avaliação do modelo, usamos conjuntos de dados de documentos jurídicos em português e inglês, representativos das diferenças culturais e linguísticas discutidas anteriormente. Esses dados serão processados para treinar o modelo com foco em identificar as nuances jurídicas, permitindo que ele compreenda as diferenças conceituais e contextuais entre termos semelhantes nas duas línguas.

\section{Métricas de avaliação}

A avaliação do sistema de tradução automática será realizada com a métrica BLEU (\textit{BiLingual Evaluation Understudy}), que mede a precisão de traduções baseando-se na similaridade entre a tradução gerada pelo sistema e traduções de referência realizadas por humanos. Esta métrica considera o alinhamento de n-gramas, o que permite avaliar a presença de termos corretos e de sequências de palavras que preservam o contexto do texto de origem. A métrica BLEU, se assemelha a análise de verdadeiro-positivo e falso-positivo, que nos permite identificar os casos em que o modelo acerta a tradução de forma contextual e aqueles em que a tradução se distancia do significado original.

Dessa forma, a métrica BLEU servirá para verificar a fidelidade do sistema em relação às traduções humanas, considerando as peculiaridades do contexto jurídico. Além disso, realizaremos análises qualitativas para identificar onde o sistema falha e onde acerta em capturar as nuances culturais e legais do texto.

\section{Treinamento e validação do modelo}

O treinamento do modelo foi dividido em etapas, incluindo a fase inicial de pré-processamento dos dados, onde os documentos jurídicos foram organizados e categorizados de acordo com seus contextos legais. A fase de treinamento do modelo utiliza \textit{embedding layers} e Modelos Transformadores (TM), que são redes neurais adequadas para capturar as relações de dados sequenciais, ou seja, contextualização.

Após o treinamento inicial, realizamos a validação cruzada utilizando amostras de textos que não foram incluídas no conjunto de treinamento. A validação permitirá identificar a capacidade do modelo de lidar com termos específicos do jargão jurídico e adaptar-se aos dois tipos de problemas previamente definidos. Eventualmente, ajustes finos serão realizados no modelo, a partir de opiniões expertos do domínio jurídico, para aprimorar sua capacidade de reconhecer e refletir nuances contextuais e culturais.

Com essa abordagem metodológica, o sistema de tradução automática se propõe a ser uma ferramenta útil para a tradução jurídica, oferecendo um suporte que considera as especificidades e os desafios da área. Os resultados obtidos e as análises dos testes serão apresentados nos próximos capítulos, abordando os pontos fortes e as limitações do sistema proposto.


% \chapter{Metodologia}


% Afim de fornecer um paradigma conceitual de um SG a Figura~\ref{fig:ModeloConceitual} demonstra os sete campos definidos pelo o Instituto Nacional de Padrões e Tecnologia (INPT)\footnote{\url{https://www.nist.gov/}}.

% \begin{figure}[h!]
%     \centering
%     \caption{Modelo Conceitual INPT}
%     \label{fig:ModeloConceitual}
%     \includegraphics[width=1\textwidth]{figs/atores.png} 

%    Fonte:  Instituto Nacional de Padrões e Tecnologia (INPT)
% \end{figure}

% \newpage

%  Geração de energia de forma não-renovável e alteração climáticas: No Brasil como mostra a  Figura \ref{fig:EMmedio}, o uso das termelétricas aumentou, e isso ocorreu devido à falta de  chuvas. 

% \begin{figure}[h!]
%     \centering
%     \caption{Aumento da Utilização de Termelétricas no Brasil}
%     \label{fig:EMmedio}
%     \includegraphics[width=0.95\textwidth]{figs/em_m_dio.png} 

%     Fonte: ...
% \end{figure}

% A Figura \ref{fig:Blecaute} apresenta alguns dos maiores blecautes ja registrados no mundo em termos de pessoas afetadas. O Brasil tem 2 entre os 8 maiores blecautes de todo o mundo~\cite{Lopes2015}. 

% \begin{figure}[h!]
% 	\centering
%     \caption{Maiores Blecautes no Mundo}
%     \label{fig:Blecaute}
% 	\includegraphics[width=0.75\textwidth]{figs/blecaute_mundo.png} 
	   
%     Fonte: \cite{Lopes2015}  
% \end{figure}



% \begin{center}
% \begin{minipage}{0.95\linewidth}
% \begin{lstlisting}[caption={Código Python responsável por desambiguar URIs}, label=python2,captionpos=t,belowcaptionskip=15pt]
% def get_dbpedia_uris(self, method: Method) -> Self:
%   dbpedia = DbpediaInterface()

%   if method == Method.RDF2VEC:
%     self.__uri_list = []

%     for entity in self.entities:
%       # Consulta para a DBpedia com o texto da entidade
%       entity_type = EntityType[entity.label_]
%       uris = dbpedia.get_possible_dbpedia_uris(
%         entity.text, 
%         entity_type
%       )
%     return self
% \end{lstlisting}
% \end{minipage}
% \end{center}



% \begin{table}[h]
% \begin{center}
%      \caption{Índices 1 para casos factíveis}
%      \begin{tabular}{| l | l | l | l |}
%      \hline Índice & LP Petro & LP 2 & Diferença\\ 
%      \hline $SES_y$& 60.5406 & 60.5492 & -0.0087\\
%      \hline $SES_u$& 1166.1464 & 1166.1464 & 2.36*$10^{-9}$ \\
%      \hline
 
%     \end{tabular}
% \label{table:indices1}
% \end{center}
% \end{table}